% ===================================================================
%  TABELUL Nr. 1 — Școala. — Liceul. — Clasa.
%
%  Ceci n'est PAS une transcription : c'est une traduction, faite en
%  2026, en roumain d'aujourd'hui. D'ou trois differences avec
%  texto/io et texto/fr, et trois seulement :
%
%   * pas de \nl ni de \cc ;
%   * pas de \begin{VUpage} ;
%   * le gras se pose sur le TERME seul, ponctuation dehors.
%
%  Les cles « %%K » sont celles de texto/io, macro pour macro pour
%  l'apparat, et les renvois \textsuperscript{(n)} visent les MEMES
%  objets de la planche.
%
%  DIXIEME COLONNE DES DIX-SEPT. Apres le basque, qui n'avait de
%  parent nulle part, une langue romane — et le repos que cela donne
%  est mesurable : le roumain postpose le complement du nom comme
%  l'ido, « ușa castelului », la porte du chateau. On attend donc zero
%  inversion, comme en catalan et en occitan.
%
%  MAIS IL EST LA PLUS EST DES LANGUES ROMANES, et le lexique le dit.
%  Le vocabulaire de base est latin comme partout — « școala »,
%  « masa », « carte » — mais tout ce qui touche a la maison, a la
%  ferme et a l'administration porte des mots slaves : « ceas »
%  l'horloge, « praf » la poussiere, « prieten » l'ami, « a citi »
%  lire. Verifier sur l'espagnol ou l'italien serait donc trompeur :
%  on verifie sur l'ido et le francais, comme partout ailleurs.
%
%  IL DECLINE ENCORE, seul parmi les langues romanes : le genitif et
%  le datif se marquent, et l'article defini se soude a la fin du mot
%  — « tabel », « tabelul ». Le gras porte donc la forme que la phrase
%  demande, article compris quand il fait corps avec le nom.
% ===================================================================

%%K t01-apar-0 apar
\VUsaut{2.0mm}
\VUcentre{12.6pt}{120}{BROȘURĂ EXPLICATIVĂ}
\VUsaut{3.2mm}
\VUcentre{8.4pt}{160}{LA}
\VUsaut{2.6mm}
\VUtitre{15.6pt}{0}{tabelele auxiliare Delmas}
\VUsaut{3.4mm}
\VUornamento{9pt}
\VUsaut{5.0mm}
\VUcentre{13.2pt}{90}{PRIMA SERIE}
\VUsaut{3.0mm}
\VUfilet{16mm}
\VUsaut{5.6mm}

%%K t01-apar-1 apar
\VUtitre{15.0pt}{60}{TABELUL Nr. 1}
\VUsaut{1.4mm}
\VUtitre{11.4pt}{0}{Școala. --- Liceul. --- Clasa.}
\VUsaut{2.2mm}
\VUfilet{20mm}
\VUsaut{6.4mm}

%%K t01-tit-1 sub
\VUpk{10.2pt}{40}{Descriere generală.}
\VUsaut{3.0mm}

%%K t01-01-1 p
1. --- În acest tabel este înfățișată o \VUgras{clasă} dintr-un \VUgras{liceu}. În
această clasă sunt opt \VUgras{mese} \textsuperscript{(18)} și opt \VUgras{bănci}
\textsuperscript{(32)}. Pe fiecare bancă stau trei elevi. \VUgras{Profesorul}
\textsuperscript{(7)} stă pe un \VUgras{scaun} \textsuperscript{(8)} în spatele
\VUgras{catedrei} sale \textsuperscript{(6)}.

%%K t01-02-1 p
2. --- La stânga catedrei se află un \VUgras{dulap} cu geam
\textsuperscript{(15)}. La dreapta --- \VUgras{tabla} \textsuperscript{(1)} cu
\VUgras{buretele} ei \textsuperscript{(2)} și cu \VUgras{creta}
\textsuperscript{(3)}.

%%K t01-02-2 p
Deasupra catedrei atârnă \VUgras{planșe} \textsuperscript{(9, 11, 12)} și o
\VUgras{hartă} \textsuperscript{(10)}.

%%K t01-03-1 p
3. --- Nu toți elevii sunt la \VUgras{locurile} lor \textsuperscript{(17)}. Ioan
\textsuperscript{(4)} stă în picioare lângă tablă; ține buretele și șterge
\VUgras{figurile geometrice}.

%%K t01-03-2 p
Petru stă în picioare lângă catedră; strânge \VUgras{lucrările scrise}
\textsuperscript{(5)} și le duce profesorului.

%%K t01-04-1 p
4. --- \VUgras{Acest} profesor este profesor de \VUgras{limbi moderne}. Îi învață
pe elevi să vorbească, să citească și să scrie în limbi străine
\VUgras{(germana, engleza, spaniola, italiana, rusa} sau \VUgras{franceza)}.

%%K t01-05-1 p
5. --- Clasa este foarte mare și foarte luminoasă. În fund se află o
\VUgras{fereastră} largă cu două \VUgras{ochiuri de aerisire}
\textsuperscript{(78)} și o \VUgras{ușă cu geam}, ca să fie aerisită și
luminată.

%%K t01-05-2 p
În această clasă nu este frig. La mijloc, ca să o încălzească, se află o
\VUgras{sobă} foarte bună \textsuperscript{(46)} cu \VUgras{burlanul} ei
\textsuperscript{(45)}.

%%K t01-05-3 p
De perete sunt bătute \VUgras{cuiere} \textsuperscript{(58)}; în ele atârnă
pălăriile și paltoanele elevilor.

%%K t01-tit-2 sub
\VUsaut{5.2mm}
\VUpk{10.2pt}{40}{Curtea.}
\VUsaut{3.6mm}

%%K t01-06-1 p
6. --- \VUgras{Ceasul} \textsuperscript{(63)} arată nouă și cinci minute.

%%K t01-06-2 p
\VUgras{Recreația} \textsuperscript{(76)} tocmai s-a sfârșit, și lecția începe din
nou. Recreația se petrece în \VUgras{curte} \textsuperscript{(75)}. Vedem câțiva
elevi care se joacă. \VUgras{Pedagogul} \textsuperscript{(74)} îi supraveghează.

%%K t01-07-1 p
7. --- În curte mai vedem și alți oameni: \VUgras{inspectorul}
\textsuperscript{(73)}, \VUgras{directorul} \textsuperscript{(71)} și
\VUgras{șeful de studii} \textsuperscript{(72)}.

%%K t01-07-2 p
Inspectorul cercetează școlile, directorul conduce așezământul, șeful de studii
veghează asupra muncii elevilor.

%%K t01-07-3 p
Al patrulea este \VUgras{portarul} \textsuperscript{(77)}. Poartă la subsuoară
catalogul, ca să însemne absențele.

%%K t01-08-1 p
8. --- În fundul curții se află \VUgras{aparatele de gimnastică}
\textsuperscript{(70)} și atelierul de \VUgras{lucru manual}
\textsuperscript{(69)}.

%%K t01-08-2 p
La dreapta tabelului se zăresc abia \VUgras{clasele} de \VUgras{muzică} și de
\VUgras{desen}.

%%K t01-noto-1 noto
\VUnotes{143.2mm}{%
(*) Și pentru \textit{nume} sfătuim să fie redate în forma latină. Numai așa se
pot ocoli nenumăratele variante. Și cum să alegem între \textit{Giovanni, Jean,
Johann, Jan, Hans, John, Ioan} și celelalte? Doar nu vom face un dicționar
deosebit al numelor. Să luăm nominativul latin și să spunem: \VUgras{Ioannes,
Ioanna; Iulius, Iulia; Iakobus; Andreas; Lukas.} (Gramatica completă).%
}

%%K t01-tit-3 sub
\VUpk{10.2pt}{40}{Descriere amănunțită.}
\VUsaut{2.4mm}

%%K t01-tit-4 sub
\VUpk{10.2pt}{40}{Elevii buni.}
\VUsaut{3.0mm}

%%K t01-09-1 p
9. --- Elevii din prima bancă la dreapta catedrei sunt \VUgras{Henric} (*),
\VUgras{Ștefan} și \VUgras{Carol}.

%%K t01-09-2 p
Henric este foarte atent și foarte silitor; îl ascultă pe profesor. Pe masă are un
\VUgras{caiet} \textsuperscript{(28)}, o \VUgras{carte} \textsuperscript{(29)}, un
\VUgras{dicționar} \textsuperscript{(30)} și o \VUgras{penară}
\textsuperscript{(31)}. Cartea lui are multe \VUgras{pagini}; în fiecare capitol
sunt câteva \VUgras{alineate}, iar în fiecare \VUgras{rând} sunt multe
\VUgras{cuvinte}.

%%K t01-09-4 p
Vecinul său Ștefan \textsuperscript{(24)} este harnic. Scrie în caiet lecția dată
pentru data viitoare. Are un \VUgras{scris} minunat. Cartea lui este închisă. I se
vede numai \VUgras{legătura} \textsuperscript{(27)}. Lângă el stă
\VUgras{compasul} său \textsuperscript{(26)}.

%%K t01-09-5 p
Carol \textsuperscript{(23)} ține cartea în mână; o deschide ca să mai repete o
dată lecția. Ca orice elev, are înaintea sa o \VUgras{călimară}
\textsuperscript{(25)}. Este ascultător.

%%K t01-10-1 p
10. --- La masa următoare stau \VUgras{Ludovic, Anton} și \VUgras{Pavel}. Ludovic
își spune \VUgras{lecția} \textsuperscript{(22)} și răspunde la \VUgras{întrebările}
profesorului. Anton \textsuperscript{(21)} este foarte neatent; uneori profesorul
chiar îl pedepsește. Pavel \textsuperscript{(20)} este un bun \VUgras{tovarăș},
prietenos și îndatoritor. Este foarte atent.

%%K t01-tit-5 sub
\VUsaut{4.2mm}
\VUpk{10.2pt}{40}{Elevii răi.}
\VUsaut{3.2mm}

%%K t01-11-1 p
11. --- În spatele lui Henric stă \VUgras{Gheorghe} \textsuperscript{(38)}. Nu este
un băiat rău, dar este \VUgras{flecar}, neatent și nu poate sta locului.
\VUgras{Ușurința} și \VUgras{neascultarea} lui aduc \VUgras{dezordine} în lecție,
și deseori merită o \VUgras{observație} aspră. \VUgras{Ciorna} lui
\textsuperscript{(35)} este murdară; face pe ea \VUgras{pete de cerneală} cu
\VUgras{penița} \textsuperscript{(34)}, din care pricină folosește necontenit
\VUgras{guma} \textsuperscript{(36)}; \VUgras{ghiozdanul} lui
\textsuperscript{(33)} este rupt, atât este de neîngrijit. De ce își aduce
\VUgras{tocul} \textsuperscript{(37)} la clasa de limbi moderne?

%%K t01-11-2 p
Vecinul său Edmund \textsuperscript{(39)} nu îl iubește. E drept că este cam
leneș, dar este tăcut și liniștit. Nu flecărește; numai că, în loc să asculte,
preferă să citească \VUgras{povești} din \VUgras{cartea sa de lectură}.

%%K t01-11-3 p
Al treilea elev din această bancă este \VUgras{Ludwig} \textsuperscript{(40)}. Este
silitor. Scrie pe o \VUgras{foaie de hârtie} albă. Înaintea sa a pus
\VUgras{sugativa} \textsuperscript{(41)}.

%%K t01-12-1 p
12. --- Elevii din a doua bancă, la stânga profesorului, sunt cu totul mici.
Primul, micul \VUgras{Emil}, încă nu știe să scrie bine; își aduce
\VUgras{tăblița} \textsuperscript{(42)}, \VUgras{condeiul} său și
\VUgras{buretele} său \textsuperscript{(43)}.

%%K t01-12-2 p
Lângă el sunt \VUgras{August} și \VUgras{Bernard} \textsuperscript{(44)}. Acesta
din urmă lucrează bine, dar are o \VUgras{înfățișare} foarte neîngrijită.

%%K t01-13-1 p
13. --- În spatele lor stau trei \VUgras{leneși}. \VUgras{Albert} nu își învață
lecțiile. \VUgras{Iacob} \textsuperscript{(47)} doarme în loc să asculte.
\VUgras{Lenea} lui îi atrage \VUgras{mustrări} și chiar \VUgras{pedepse}. În
sfârșit, \VUgras{Cristian} \textsuperscript{(48)} este prea ușuratic. Se
\VUgras{poartă} rău și nu se străduiește să se îndrepte.

%%K t01-14-1 p
14. --- Vecinii lor, dimpotrivă, se poartă bine. \VUgras{Ernest}
\textsuperscript{(51)} iubește ordinea și exactitatea; se folosește întotdeauna de
\VUgras{rigla} sa \textsuperscript{(50)} și de \VUgras{creionul} său
\textsuperscript{(49)}. Este \VUgras{elev extern}.

%%K t01-14-2 p
\VUgras{Francisc} \textsuperscript{(52)} este \VUgras{intern}; poartă uniformă,
mănâncă și doarme în așezământ. Este un bun \VUgras{tovarăș}.

%%K t01-14-3 p
Mauriciu își ascute \VUgras{creionul} cu \VUgras{briceagul} său
\textsuperscript{(56)}; este foarte chibzuit.

%%K t01-14-4 p
Își aduce toate cărțile, \VUgras{gramatica} sa \textsuperscript{(55)},
\VUgras{carnetul} său \textsuperscript{(54)} și chiar \VUgras{blocnotesul}
\textsuperscript{(53)}. \VUgras{Servieta} lui \textsuperscript{(57)} stă pe
podea, în spatele lui.

%%K t01-14-5 p
Cele două mese din fundul clasei sunt ocupate de \VUgras{elevii buni}
\textsuperscript{(19)}.

%%K t01-tit-6 sub
\VUsaut{4.6mm}
\VUpk{10.2pt}{40}{Desenul și muzica.}
\VUsaut{3.4mm}

%%K t01-15-1 p
15. --- În \VUgras{clasa de desen} atârnă pe perete \VUgras{modele} frumoase, iar
din tavan atârnă o \VUgras{lampă} \textsuperscript{(62)} cu \VUgras{abajurul} ei.
\VUgras{Profesorul} \textsuperscript{(60)} este foarte răbdător; îndreaptă
\VUgras{desenele} \textsuperscript{(61)} elevilor și îi învață să deseneze
\VUgras{bustul} \textsuperscript{(59)} după natură.

%%K t01-16-1 p
16. --- \VUgras{Clasa de muzică} pare foarte plăcută. Acolo se învață să se
citească \VUgras{notele}, să se cânte \VUgras{treptele gamei}
\textsuperscript{(67)}, scrise pe \VUgras{portativ} (do, re, mi, fa, sol, la,
si\,=\,C. D. E. F. G. A. B.), să se cunoască \VUgras{cheile} și să se cânte
\VUgras{melodii} încântătoare. Notele se pun pe \VUgras{pupitru}
\textsuperscript{(65)}. Unul dintre elevi \VUgras{cântă la pian}
\textsuperscript{(64)}, iar \VUgras{profesorul de muzică} \textsuperscript{(66)}
\VUgras{bate tactul} \textsuperscript{(68)}.

%%K t01-17-1 p
17. --- Se zărește abia un colț al atelierului de \VUgras{lucru manual}
\textsuperscript{(69)}, unde copiii pot învăța să rindeluiască lemnul și să
pilească fierul. Nu în orice școală se află așa ceva.

%%K t01-tit-7 sub
\VUsaut{5.0mm}
\VUpk{10.2pt}{40}{Clasa de științe naturale.}
\VUsaut{3.6mm}

%%K t01-18-1 p
18. --- Profesorul de matematică îi învață pe elevi \VUgras{geometria,}
\VUgras{aritmetica, socotitul} și \VUgras{sistemul metric}. Pe tabel vedem
figurile geometrice de căpetenie: \VUgras{cercul} \textsuperscript{(a)},
\VUgras{pătratul} \textsuperscript{(b)}, \VUgras{triunghiul}
\textsuperscript{(c)}, \VUgras{linia} \textsuperscript{(d)}.

%%K t01-18-2 p
Vedem de asemenea o \VUgras{operație}; nu este \VUgras{adunare}, nici
\VUgras{scădere}, nici \VUgras{înmulțire}: este \VUgras{împărțire}
\textsuperscript{(e)}.

%%K t01-19-1 p
19. --- Pe tabelul sistemului metric vedem: \VUgras{metrul}
\textsuperscript{(a)}, \VUgras{balanța} \textsuperscript{(b)} și
\VUgras{greutățile} ei, \VUgras{litrul} \textsuperscript{(e)},
\VUgras{decalitrul} \textsuperscript{(f)}, \VUgras{decilitrul} și \VUgras{metrul
cub} \textsuperscript{(d)}.

%%K t01-19-2 p
Elevii mai învață și \VUgras{științele naturale} \textsuperscript{(11)} ---
zoologia \textsuperscript{(a)}, botanica \textsuperscript{(b)}, geologia
\textsuperscript{(c)} --- și \VUgras{istoria} \textsuperscript{(12)}.

%%K t01-19-3 p
În dulap se află aparate pentru \VUgras{fizică} \textsuperscript{(15)} și pentru
\VUgras{chimie} \textsuperscript{(16)}.

%%K t01-19-4 p
Pe dulap se află: \VUgras{sfera} \textsuperscript{(14)}, \VUgras{conul} și
\VUgras{globul} \textsuperscript{(13)}.

%%K t01-19-5 p
Deasupra tabelelor atârnă o \VUgras{hartă} care înfățișează cele cinci părți ale
lumii: \VUgras{Europa} \textsuperscript{(g)}, \VUgras{Asia}
\textsuperscript{(e)}, \VUgras{Africa} \textsuperscript{(f)}, \VUgras{America}
\textsuperscript{(ab)} și \VUgras{Oceania} \textsuperscript{(h)}, și cele două
mari oceane --- \VUgras{Pacificul} \textsuperscript{(c)} și \VUgras{Atlanticul}
\textsuperscript{(d)}.
\VUsaut{6.0mm}
\VUfilet{22mm}
